\chapter{绪论}

在平台化传播不断重组文娱产业的背景下，艺人的职业发展越来越像一个需要持续调度资源、控制节奏并
维护多方关系的复杂项目。传统媒体时代，杂志封面、电视剧角色与综艺节目往往构成少数关键入口；
而在社交媒体时代，内容更新频率、与受众的互动密度以及对自我叙事的控制能力，逐渐成为影响职业
韧性的核心变量。以公开人物案例来观察这一变化，有助于把抽象的“转型”拆解为可比较、可复用的
分析步骤。

\section{研究背景}

佐佐木希的职业经历兼具高识别度、跨阶段明显与公开资料丰富三项特点。她从平面模特进入公众视野，
随后拓展到电视剧、电影、综艺与广告等多种媒介形态，并在争议事件后通过社交平台与生活化内容
逐步重建公共形象。这个过程既包含传统媒体中的角色扩展，也包含平台时代的内容再组织，因而很适合
作为毕业论文模板的公开演示案例。

对 `thesis-nju-master` 而言，选择这样一个主题还有额外好处。它可以帮助用户在不暴露真实研究对象、
样本数据与私人信息的前提下，直接看到模板在封面、摘要、目录、章节、图表、表格、参考文献和附录
等页面上的完整表现。也就是说，这份示例正文既是一篇可以阅读的短论文，也是项目级模板说明书。

\section{研究问题}

本文围绕以下三个问题展开：

\begin{enumerate}
  \item 佐佐木希的跨媒介职业路径可以划分为哪些具有辨识度的阶段？
  \item 在不同阶段中，哪些内容运营动作帮助她维持或修复公众可见性？
  \item 如何把上述分析写成一份既像真实论文、又能帮助新用户快速上手模板的公开示例？
\end{enumerate}

\section{研究思路与模板展示策略}

本文采用“公开案例梳理 + 模板功能展示”相结合的写法。研究思路上，先依据公开时间线划分职业阶段，
再从媒介入口、内容节奏与形象韧性三个角度总结各阶段的关键特征；展示策略上，则把摘要、章节、
图表、表格、参考文献、致谢与附录全部保留下来，让用户打开 `main.tex` 即可看到一份结构完整的
示例论文，而不需要先理解复杂的隐藏工作流。

\section{章节安排}

第一章说明问题背景、研究问题与模板展示策略。第二章进入案例分析，并结合图、表与附录示例展示
`thesis-nju-master` 在南京大学工程管理硕士论文场景下的基本排版能力。这样安排的目的，是让使用者
既能把它当作一篇可读的公开案例，也能把它当作修改自己论文时的直接参考。
